\subsection{FPTRANSA Common Model}\label{section:fptransa-common-model}
For an approximate transcendental instruction, $x'$ is the source after \texttt{FSTATUS.DAZ}, $y=f(x')$ is the
exact real reference value named by its accuracy contract, and $a$ is the approximate result before
\texttt{FSTATUS.FTZ}. For a format with precision $p$ and minimum normal exponent $e_{\min}$, define

\[
  \operatorname{ulp}_F(y) =
  \begin{cases}
    2^{\max(\lfloor\log_2 |y|\rfloor-p+1,\ e_{\min}-p+1)}, & y\ne0,\\
    2^{e_{\min}-p+1}, & y=0.
  \end{cases}
\]

The measured error is $|a-y|/\operatorname{ulp}_F(y)$. The bound applies only to finite reference values within
the selected format's finite range; NaN, infinity, and overflow follow the instruction's special-value rules.
Every present contract guarantees at most 4 ULP for both S and D, corresponding to at least 21 and 50 worst-case
relative precision bits respectively for nonzero normal results. CPUID may advertise a smaller certified bound.

FPTRANSA applies DAZ before forming the exact reference input and measures its advertised ULP error before FTZ.
It ignores \texttt{FSTATUS.RM} and formats its implementation-defined approximation with nearest-even rounding.
The approximation is deterministic for the same implementation, input, format, and relevant FSTATUS mode, but need
not be bit-identical across implementations.

A signaling NaN accrues NV. NaN results are quieted and then processed by \texttt{FSTATUS.DN}.
\texttt{FSTATUS.FTZ} replaces a subnormal result with the required signed zero and accrues UF. Each instruction
accrues NV and DZ according to its special-value rules, and accrues OF and UF when required by its exact result;
NX remains unchanged and cannot cause a trap. Overflow means that the exact result magnitude exceeds the selected
format's maximum finite value. Underflow means that the exact nonzero result magnitude is below the selected format's
minimum normal value. If any accrued exception is enabled as a trap, the instruction raises
\texttt{FLOATING\_POINT\_FAULT} before writing any destination; otherwise it commits its result or results and
accrued flags together.

\clearpage
